% CPSY 1291 -- Recitation handout 14
% Compile:  latexmk -lualatex handout-14-project-clinic-ii.tex
\documentclass[11pt]{article}
\usepackage{cpsy1291-handout}

\begin{document}

\handouthead{14}{Project clinic II}
  {Figures that are instruments, a report that is honest, and a talk that lands}
  {Week 14 \ $\cdot$\ last class meeting Thu 12/10; presentations Mon 12/21, 9:00\,AM}

\section*{What this session is for}
\addcontentsline{toc}{section}{What this session is for}

By now you have results, or you have discovered that you will not get the ones
you wanted. Both are fine, and the syllabus is explicit that a negative result
clearly analyzed is a perfectly good project. What is graded is the quality of
the question, the soundness of the method, and how honestly you interpret what
you found.

This handout is about turning what you have into three things: figures that
carry an argument, a report that survives a skeptical reader, and a talk that
lands in the time you are given.

\begin{keybox}
\ask{The one idea.} \textbf{The claim comes first, and everything else is
evidence for it.} Write your one-sentence result before you make a single
figure. Then every figure either supports that sentence or comes out. Projects
that assemble figures first and look for a story afterwards produce talks that
are a tour of what was done rather than an argument about what is true.
\end{keybox}

\section{Figures as instruments}

\subsection{One figure, one claim}

Each figure should be describable in one sentence that is a claim, not a
description. ``Accuracy by condition and model'' is a description. ``The texture
model loses more accuracy under shape distortion than the shape model'' is a
claim, and it tells you what to plot and what to leave out.

If a figure needs two sentences, it is two figures. If it needs none, it is not a
figure.

\subsection{The checklist}

\begin{center}
\begin{tabular}{@{}ll@{}}
\toprule
Every figure & Why \\
\midrule
axis labels with \textbf{units} & an unlabeled axis is uninterpretable \\
\textbf{error bars}, and a caption saying which & SD and SE answer different questions \\
a \textbf{chance line} where one exists & a reader cannot infer it \\
$n$ stated --- of \emph{what} & runs? stimuli? subjects? \\
readable at the size it will be seen & shrink it to slide size and check \\
distinguishable without color & about 8\% of men are color-blind \\
\bottomrule
\end{tabular}
\end{center}

\begin{keybox}
\ask{The error-bar rule, once more.} Label them. \emph{Standard deviation} says
how much a single run varies; \emph{standard error} says how well you know the
mean. An unlabeled error bar cannot be read, and a figure whose error bars are
unlabeled is asking the reader to trust you rather than to check you.
\end{keybox}

\subsection{Three plots that are almost always worth making}

\textbf{The distribution, not just the mean.} A bar chart of two means hides
whether the distributions overlap completely. Plot the individual points over the
bars --- with ten runs there is no reason not to.

\textbf{The thing that would falsify you.} If your claim is that model A beats
model B, plot the per-item differences, not just the two averages. If A wins on
51\% of items by a hair, that is a different claim from A winning on 90\%.

\textbf{The failure cases.} Show the items your model got wrong. It is the most
informative figure in most projects and it is almost always missing, because
nobody wants to put their failures on a slide. It is also what makes a talk
credible.

\subsection{Two things not to do}

\textbf{Do not truncate a $y$-axis} to make a difference look big. If the
difference needs a truncated axis, the difference is small, and saying so is
better than being caught.

\textbf{Do not use a rainbow colormap} for continuous data. It creates
boundaries that are not in the data. \texttt{viridis} is the default for a
reason.

\section{The report}

\subsection{A structure that always works}

\begin{center}
\begin{tabular}{@{}ll@{}}
\toprule
Section & The one job it does \\
\midrule
Question & what you asked, and why it is worth asking \\
Method & enough that someone else could redo it \\
Results & what happened --- figures, numbers, no interpretation \\
Interpretation & what it means, and what it does not \\
Limitations & what would change your mind \\
\bottomrule
\end{tabular}
\end{center}

The split between Results and Interpretation is the one that matters. Keeping
them separate is what lets a reader disagree with your interpretation while
accepting your results, which is how science is supposed to work --- and it is
what makes a report readable by someone who does not already agree with you.

\subsection{Limitations is a strength section}

Students write limitations defensively, as damage control. Written well it is
the section that demonstrates you understand your own work.

Say what would change your mind. ``We used one architecture, so we cannot
separate an effect of depth from an effect of this particular network'' is a
strong sentence: it names the confound, bounds the claim, and tells the reader
exactly which experiment would settle it. ``More data would help'' is not a
limitation; it is a truism.

\subsection{The AI disclosure}

The syllabus permits AI tools for the project and requires a short disclosure
saying which tools you used and how they contributed. Write it plainly and
specifically --- ``we used an LLM to draft the plotting code and to debug the data
loader'' --- rather than vaguely. It costs two lines, it is a course requirement,
and vagueness reads worse than the use itself.

\section{The talk}

\subsection{The arc}

\begin{center}
\begin{tabular}{@{}ll@{}}
\toprule
Slide(s) & Job \\
\midrule
1 & the question, in one sentence, and why anyone should care \\
2 & what you did, in one picture --- not in bullets \\
3--5 & the result: one figure per slide, one claim per figure \\
6 & the honest caveat: what would change your mind \\
7 & what you would do next \\
\bottomrule
\end{tabular}
\end{center}

\begin{keybox}
\ask{Slide 2 is the one people get wrong.} Your method is a pipeline, and a
pipeline is a picture --- boxes and arrows, from data to number. Three bullet
points describing the same pipeline take longer to say and are harder to follow.
Draw it.
\end{keybox}

\subsection{Rehearsal, which is not optional}

\begin{enumerate}
  \item \textbf{Out loud, with a timer, once.} Silent reading takes roughly half
        the time speaking does, so a deck that feels right unspoken is about
        twice too long.
  \item \textbf{Decide who says what.} In a group talk, silent handovers are the
        commonest thing to go wrong. Write the handover into the notes.
  \item \textbf{Cut to fit.} If you are over, cut a \emph{result}, not the caveat
        and not the question. A talk that covers less and lands is better than
        one that covers everything at speed.
\end{enumerate}

\subsection{Questions you will be asked}

\begin{itemize}
  \item \emph{What is chance here?} --- have the number, not a gesture.
  \item \emph{How many seeds / subjects / stimuli?} --- have the $n$, and say what
        it is the $n$ \emph{of}.
  \item \emph{What does a trivial baseline get?} --- if you do not have one, say
        so plainly; do not improvise.
  \item \emph{Would this hold for a different dataset or architecture?} --- the
        honest answer is usually ``we do not know, and here is why we think it
        might not''.
  \item \emph{Why this metric?} --- have a reason, even a modest one.
\end{itemize}

\begin{keybox}
\ask{``We don't know'' is a good answer.} Said once, with a reason, it reads as
command of the material. What reads badly is a confident answer to a question you
did not investigate --- and the room can tell the difference.
\end{keybox}

\section{If the result did not come out}

This happens, and it is survivable. The syllabus grades the question, the method
and the honesty of the interpretation --- not whether the model worked.

\begin{enumerate}
  \item \textbf{Say what you expected and what you got.} Both, plainly, early.
  \item \textbf{Show that the pipeline works.} A positive control --- a case where
        your method \emph{does} detect a known effect --- turns ``our analysis
        found nothing'' into ``there is nothing there to find, and here is the
        evidence that we could have found it''. Those are very different claims
        and the second is a result.
  \item \textbf{Bound it.} How large an effect could you have detected? A null
        with a stated sensitivity is informative; a null without one is silence.
  \item \textbf{Say what you would do differently.} This is the part that shows
        you learned something, which is what the exercise was for.
\end{enumerate}

\section{The last checklist}

Run this the night before.

\begin{enumerate}
  \item Can you state the result in one sentence, without hedging?
  \item Does every figure have labeled axes, labeled error bars, and a stated $n$?
  \item Is chance marked on every figure that has one?
  \item Does the talk have a positive control or a trivial baseline somewhere?
  \item Have you said out loud what would change your mind?
  \item Did you rehearse with a timer, out loud, at least once?
  \item Does the report separate what happened from what you think it means?
  \item Is the AI-use disclosure written?
  \item Can someone else run your pipeline from your loader?
  \item Have you cut the slide you are keeping only because the work was hard?
\end{enumerate}

Item 10 catches something in almost every project.

\end{document}
